% figures/w04-ssa-circle.svg — 왜 선수각 오차를 ssa 로 감아야 하는가.
%
%   bash _tools/tikz2svg.sh figures/src/w04-ssa-circle.tex figures/w04-ssa-circle.svg
%
% 상황: psi = 170 deg 를 유지하는 배에 psi_d = -170 deg 를 명령한다.
%   뺄셈 그대로       e = -170 - 170 = -340 deg  -> 좌현으로 340 도
%   감아서            ssa(e) = +20 deg           -> 우현으로 20 도
%
% 각 계산 (NED: 북에서 시계방향). 화면각 = 90 - 나침반각.
%   psi   = 170 -> 화면 -80
%   psi_d = 190 -> 화면 -100  (= 260)
% 짧은 호: 화면 -80 -> -100   (시계방향 20 도)
% 긴 호  : 화면 -80 -> 260    (반시계 340 도)
\documentclass[border=5pt]{standalone}
% gnc-style.tex — 이 볼트의 TikZ 그림이 공유하는 스타일.
%
% 저널 그림 규약을 스타일 하나로 굳혀 둔다. 그림마다 선 굵기나 화살촉을
% 다시 정하지 않는다 — 그것이 그림이 제각각으로 보이는 원인이다.
%
% 쓰는 법:  % gnc-style.tex — 이 볼트의 TikZ 그림이 공유하는 스타일.
%
% 저널 그림 규약을 스타일 하나로 굳혀 둔다. 그림마다 선 굵기나 화살촉을
% 다시 정하지 않는다 — 그것이 그림이 제각각으로 보이는 원인이다.
%
% 쓰는 법:  % gnc-style.tex — 이 볼트의 TikZ 그림이 공유하는 스타일.
%
% 저널 그림 규약을 스타일 하나로 굳혀 둔다. 그림마다 선 굵기나 화살촉을
% 다시 정하지 않는다 — 그것이 그림이 제각각으로 보이는 원인이다.
%
% 쓰는 법:  \input{gnc-style}  (figures/src/ 안에서)

\usepackage{tikz}
\usepackage{amsmath}
\usepackage{xcolor}
\usetikzlibrary{arrows.meta, positioning, calc, fit, backgrounds, decorations.pathmorphing}

% ---- 색 --------------------------------------------------------------------
% 색은 강조 하나에만 쓴다. 나머지는 검정.
\definecolor{ink}{HTML}{1A1A1A}
\definecolor{accent}{HTML}{7C3AED}
\definecolor{warn}{HTML}{B91C1C}
\definecolor{soft}{HTML}{666666}

% ---- 선과 화살촉 -----------------------------------------------------------
\tikzset{
  % 신호선. 화살촉은 작고 뾰족한 것 하나뿐이다.
  sig/.style   = {ink, line width=0.5pt, -{Stealth[length=4pt, width=3pt]}},
  wire/.style  = {ink, line width=0.5pt},
  % 강조 경로 — 파선으로 구분한다. 흑백 인쇄에서도 살아야 하기 때문이다.
  asig/.style  = {accent, line width=0.5pt, densely dashed,
                  -{Stealth[length=4pt, width=3pt]}},
  awire/.style = {accent, line width=0.5pt, densely dashed},
  % 블록. 흰 바탕, 직각, 검은 테두리.
  blk/.style   = {draw=ink, line width=0.55pt, fill=white,
                  minimum width=17mm, minimum height=8mm, inner sep=2pt, font=\small},
  % 합산점
  sum/.style   = {draw=ink, line width=0.55pt, fill=white, circle,
                  inner sep=0pt, minimum size=4.6mm, font=\scriptsize},
  asum/.style  = {draw=accent, line width=0.55pt, fill=white, circle,
                  inner sep=0pt, minimum size=4.6mm, font=\scriptsize},
  % 분기점
  dot/.style   = {ink, fill=ink, circle, inner sep=0pt, minimum size=1.5mm},
  % 신호 이름 — 선 위에 작게
  sname/.style = {font=\small, inner sep=1.5pt},
  % 그림 안의 짧은 설명. 그림 안의 글은 최소로 둔다
  note/.style  = {font=\scriptsize, soft, inner sep=1.5pt},
  anote/.style = {font=\scriptsize, accent, inner sep=1.5pt},
  % 축
  axis/.style  = {ink, line width=0.5pt},
  tick/.style  = {font=\scriptsize, soft},
}

% 캡션 한 줄 — 그림 아래 가는 선과 함께
\newcommand{\gnccaption}[2]{%
  \node[anchor=north west, text width=#1, align=left, font=\scriptsize, ink]
        at (current bounding box.south west) {\vspace{2pt}#2};}
  (figures/src/ 안에서)

\usepackage{tikz}
\usepackage{amsmath}
\usepackage{xcolor}
\usetikzlibrary{arrows.meta, positioning, calc, fit, backgrounds, decorations.pathmorphing}

% ---- 색 --------------------------------------------------------------------
% 색은 강조 하나에만 쓴다. 나머지는 검정.
\definecolor{ink}{HTML}{1A1A1A}
\definecolor{accent}{HTML}{7C3AED}
\definecolor{warn}{HTML}{B91C1C}
\definecolor{soft}{HTML}{666666}

% ---- 선과 화살촉 -----------------------------------------------------------
\tikzset{
  % 신호선. 화살촉은 작고 뾰족한 것 하나뿐이다.
  sig/.style   = {ink, line width=0.5pt, -{Stealth[length=4pt, width=3pt]}},
  wire/.style  = {ink, line width=0.5pt},
  % 강조 경로 — 파선으로 구분한다. 흑백 인쇄에서도 살아야 하기 때문이다.
  asig/.style  = {accent, line width=0.5pt, densely dashed,
                  -{Stealth[length=4pt, width=3pt]}},
  awire/.style = {accent, line width=0.5pt, densely dashed},
  % 블록. 흰 바탕, 직각, 검은 테두리.
  blk/.style   = {draw=ink, line width=0.55pt, fill=white,
                  minimum width=17mm, minimum height=8mm, inner sep=2pt, font=\small},
  % 합산점
  sum/.style   = {draw=ink, line width=0.55pt, fill=white, circle,
                  inner sep=0pt, minimum size=4.6mm, font=\scriptsize},
  asum/.style  = {draw=accent, line width=0.55pt, fill=white, circle,
                  inner sep=0pt, minimum size=4.6mm, font=\scriptsize},
  % 분기점
  dot/.style   = {ink, fill=ink, circle, inner sep=0pt, minimum size=1.5mm},
  % 신호 이름 — 선 위에 작게
  sname/.style = {font=\small, inner sep=1.5pt},
  % 그림 안의 짧은 설명. 그림 안의 글은 최소로 둔다
  note/.style  = {font=\scriptsize, soft, inner sep=1.5pt},
  anote/.style = {font=\scriptsize, accent, inner sep=1.5pt},
  % 축
  axis/.style  = {ink, line width=0.5pt},
  tick/.style  = {font=\scriptsize, soft},
}

% 캡션 한 줄 — 그림 아래 가는 선과 함께
\newcommand{\gnccaption}[2]{%
  \node[anchor=north west, text width=#1, align=left, font=\scriptsize, ink]
        at (current bounding box.south west) {\vspace{2pt}#2};}
  (figures/src/ 안에서)

\usepackage{tikz}
\usepackage{amsmath}
\usepackage{xcolor}
\usetikzlibrary{arrows.meta, positioning, calc, fit, backgrounds, decorations.pathmorphing}

% ---- 색 --------------------------------------------------------------------
% 색은 강조 하나에만 쓴다. 나머지는 검정.
\definecolor{ink}{HTML}{1A1A1A}
\definecolor{accent}{HTML}{7C3AED}
\definecolor{warn}{HTML}{B91C1C}
\definecolor{soft}{HTML}{666666}

% ---- 선과 화살촉 -----------------------------------------------------------
\tikzset{
  % 신호선. 화살촉은 작고 뾰족한 것 하나뿐이다.
  sig/.style   = {ink, line width=0.5pt, -{Stealth[length=4pt, width=3pt]}},
  wire/.style  = {ink, line width=0.5pt},
  % 강조 경로 — 파선으로 구분한다. 흑백 인쇄에서도 살아야 하기 때문이다.
  asig/.style  = {accent, line width=0.5pt, densely dashed,
                  -{Stealth[length=4pt, width=3pt]}},
  awire/.style = {accent, line width=0.5pt, densely dashed},
  % 블록. 흰 바탕, 직각, 검은 테두리.
  blk/.style   = {draw=ink, line width=0.55pt, fill=white,
                  minimum width=17mm, minimum height=8mm, inner sep=2pt, font=\small},
  % 합산점
  sum/.style   = {draw=ink, line width=0.55pt, fill=white, circle,
                  inner sep=0pt, minimum size=4.6mm, font=\scriptsize},
  asum/.style  = {draw=accent, line width=0.55pt, fill=white, circle,
                  inner sep=0pt, minimum size=4.6mm, font=\scriptsize},
  % 분기점
  dot/.style   = {ink, fill=ink, circle, inner sep=0pt, minimum size=1.5mm},
  % 신호 이름 — 선 위에 작게
  sname/.style = {font=\small, inner sep=1.5pt},
  % 그림 안의 짧은 설명. 그림 안의 글은 최소로 둔다
  note/.style  = {font=\scriptsize, soft, inner sep=1.5pt},
  anote/.style = {font=\scriptsize, accent, inner sep=1.5pt},
  % 축
  axis/.style  = {ink, line width=0.5pt},
  tick/.style  = {font=\scriptsize, soft},
}

% 캡션 한 줄 — 그림 아래 가는 선과 함께
\newcommand{\gnccaption}[2]{%
  \node[anchor=north west, text width=#1, align=left, font=\scriptsize, ink]
        at (current bounding box.south west) {\vspace{2pt}#2};}


\begin{document}
\begin{tikzpicture}[x=1mm, y=1mm]

\def\R{28}          % 나침반 원 반지름 / compass radius
\def\PS{-80}        % psi   = 170 deg
\def\PD{-100}       % psi_d = 190 deg = -170 deg

% ---- 제목 ------------------------------------------------------------------
\node[font=\small\bfseries, anchor=north west] at (-46,56)
     {Why the heading error is wrapped};
\node[anchor=north west, align=left, text width=104mm, font=\scriptsize, soft] at (-46,50.5) {%
  The vessel holds $\psi = 170^\circ$ and is commanded to $\psi_d = -170^\circ$.
  The two headings are $20^\circ$ apart, on either side of the $\pm180^\circ$ seam.};

% ---- 나침반 원과 눈금 ------------------------------------------------------
\draw[soft, line width=0.4pt] (0,0) circle (\R);
\foreach \c in {0,30,...,330} {
  \draw[soft, line width=0.35pt] ({90-\c}:\R) -- ({90-\c}:{\R-1.6});
}
\foreach \c/\lab in {0/N, 90/E, 270/W} {
  \node[font=\scriptsize, soft] at ({90-\c}:{\R+3.4}) {\lab};
}
\node[font=\scriptsize, soft, anchor=south] at ({90-45}:{\R+2.2}) {$45^\circ$};
\node[font=\scriptsize, soft, anchor=north] at ({90-135}:{\R+2.2}) {$135^\circ$};
\node[font=\scriptsize, soft, anchor=north] at ({90-225}:{\R+2.2}) {$225^\circ$};
\node[font=\scriptsize, soft, anchor=south] at ({90-315}:{\R+2.2}) {$315^\circ$};

% ---- 이음매 ($\pm 180$), 남쪽 ----------------------------------------------
\draw[warn, line width=0.45pt, densely dotted] (0,0) -- ({90-180}:{\R+14});
\node[font=\scriptsize, warn, anchor=north east, align=right] at (-4,{-\R-15})
     {S --- the seam\\$+180^\circ \equiv -180^\circ$};

% ---- 두 선수각 -------------------------------------------------------------
\draw[axis, -{Stealth[length=4.5pt,width=3.4pt]}] (0,0) -- (\PS:\R);
\draw[accent, line width=0.6pt, -{Stealth[length=4.5pt,width=3.4pt]}] (0,0) -- (\PD:\R);
\draw[soft, line width=0.3pt] ({\PS}:{\R*0.78}) -- (22,-15);
\node[font=\scriptsize, ink, anchor=west, align=left] at (22,-15)
     {$\psi = 170^\circ$\\holding};
\draw[soft, line width=0.3pt] ({\PD}:{\R*0.78}) -- (-22,-15);
\node[font=\scriptsize, accent, anchor=east, align=right] at (-22,-15)
     {$\psi_d = -170^\circ$\\commanded};

% ---- 짧은 호: ssa 를 거친 명령. 원 바깥에 그려 크게 보이게 ------------------
\draw[accent, line width=0.8pt, -{Stealth[length=4.5pt,width=3.4pt]}]
      ({\PS}:{\R+7}) arc[start angle=\PS, end angle=\PD, radius={\R+7}];
\draw[accent, line width=0.35pt] ({-88}:{\R+7}) -- (14,{-\R-14});
\node[font=\scriptsize, accent, anchor=north west, align=left] at (13,{-\R-14})
     {with \texttt{ssa}:\\$+20^\circ$ to starboard};

% ---- 긴 호: 뺄셈 그대로, 원 안쪽 ------------------------------------------
\draw[warn, line width=0.8pt, -{Stealth[length=4.5pt,width=3.4pt]}]
      ({\PS}:{\R*0.72}) arc[start angle=\PS, end angle=260, radius={\R*0.72}];
\node[font=\scriptsize, warn, anchor=north, align=center] at (0,{\R*0.72-2})
     {without \texttt{ssa}:\\$-340^\circ$ to port};

% ---- 계산 상자 -------------------------------------------------------------
\node[draw=soft, line width=0.4pt, fill=white, anchor=north west, align=left,
      font=\scriptsize, inner sep=3pt, text width=44mm] at (36,26) {%
  $e = \psi_d - \psi = -170 - 170 = -340^\circ$\\[3pt]
  $\mathrm{ssa}(e) = \mathrm{atan2}(\sin e, \cos e)$\\[1pt]
  \phantom{$\mathrm{ssa}(e)$} $= +20^\circ$\\[4pt]
  Both angles name the same turn target.\\
  Only the second one is the short way round.};

\node[anchor=north west, align=left, text width=44mm, font=\scriptsize, soft] at (36,-2) {%
  Measured in Experiment 4-6: the vessel turns $20.0^\circ$ with the wrap and
  $340.0^\circ$ without it --- seventeen times further, for the same command.};

\end{tikzpicture}
\end{document}
